\chapter{Molecular Films and Their Resonant Optical Response}\label{chap:molecular-films}

This chapter sets out the materials science and the optical physics on which the later chapters rest. It begins with the device properties that molecular packing controls and with the growth of vapor-deposited films of the two molecules studied here, ZnPc and Alq\textsubscript{3}. It then treats a single C 1s excitation. The orbitals of a conjugated bond fix the direction of its transition dipole and with it the contrast the resonance produces, a driven oscillator built on that dipole yields a causal pair of optical constants, and DFT calculations reduced by clustering supply the oscillators of a whole molecule. The chapter closes by projecting the refined molecular tensor into the frame of a film and testing a film of uniform tilt against angle-resolved NEXAFS of ZnPc.

\section{Packing and the Emergent Properties of Molecular Films}\label{sec:packing-and-the-emergent-properties-of-molecular-films}

\begin{figure}[htp]
  \centering
  \includegraphics[width=\linewidth,height=0.75\textheight,keepaspectratio]{hero/hero_fig2-01_packing}
  \caption[Molecular orientation sets the function of organic devices.]{Molecular orientation sets the function of organic devices. Molecules are drawn as slabs seen edge-on and transition dipoles as red arrows. (a) Organic transistor, in which molecules standing edge-on between source (S) and drain (D) stack their π systems along the channel and carry charge in the plane of the film. (b) Bilayer organic photovoltaic with face-on donor and acceptor layers between ITO and Al. Light entering through the glass is absorbed in the donor, and charge separates at the heterointerface, the electron passing into the acceptor and the hole returning through the donor. (c) Organic light-emitting diode. Emitters lying flat carry in-plane dipoles whose light escapes through the ITO (left), while upright emitters radiate along the film and their light remains trapped in the stack (right). (d) Organic spin valve, in which the first molecular layer on the Co electrode (red) hybridizes with the metal, shifting and broadening its HOMO and LUMO relative to the bulk film. The property set by orientation and a representative device are named beneath each panel.}
  \label{fig:2.1}
\end{figure}

Conjugated molecules interact with their neighbors through the overlap of their π orbitals, and the overlap depends on how the molecules are arranged. The electronic coupling between two molecules falls off steeply with their separation and with the angle between their π systems, so a film of identical molecules can conduct well along one direction and poorly along another.\autocite{bredas2004charge} In a planar molecule such as ZnPc the π orbitals stand normal to the molecular plane, so the strongest coupling runs along the stacking direction between cofacial neighbors and the coupling between neighbors that meet edge to edge is weak. The optical transitions of such a molecule are anisotropic as well, since its π to π* transitions have their dipoles in the molecular plane. The molecule absorbs and emits light polarized in that plane and absorbs light polarized along its normal only weakly, the reverse of the C 1s to π* excitations treated later in this chapter, whose dipoles stand normal to the plane. Nearly every property of a molecular film that matters in a device is anisotropic in this way, and the direction in which each property is large is set by the orientation of the molecules relative to the electrodes and interfaces around them. A molecule lying face-on, with its plane parallel to the substrate, stacks its π systems through the thickness of the film and lays its optical transition dipoles in the plane of the film. A molecule standing edge-on, with its plane perpendicular to the substrate, stacks its π systems within the plane of the film and turns part of its optical response toward the film normal. \Cref{fig:2.1} shows four classes of device in which this dependence is decisive, and no single orientation serves all four.

In a field-effect transistor charge moves between source and drain electrodes that lie in the plane of the film (\cref{fig:2.1}a). Molecules standing edge-on to the dielectric stack their π systems along the channel, and polymer films with this orientation reach mobilities about two orders of magnitude higher than films of the same polymer lying flat.\autocite{sirringhaus1999two} Face-on molecules in the same channel pass charge between neighbors only through the weak coupling of molecules that meet edge to edge, or across the boundaries between domains of different orientation. The current flows in the first few molecular layers above the dielectric\autocite{dodabalapur1995organic}, so the orientation at a buried interface sets device performance.\autocite{dindault2022correlating} The gate field draws the accumulated charge against the dielectric and holds it there, so a film whose bulk is well ordered can still perform poorly when its first layer adopts a different orientation, packs less densely, or breaks into small domains.

A photovoltaic cell requires the opposite arrangement (\cref{fig:2.1}b). Light enters through a transparent electrode, excitons form in the donor, and charge separates at the donor-acceptor heterointerface before it travels through the thickness of the film to the electrodes. Face-on molecules stack their π systems along that path and lay their transition dipoles in the plane of the film, parallel to the electric field of light arriving along the normal, so they absorb that light more strongly than edge-on molecules, which present only part of their transition dipole to the field. Charge separation and recombination both proceed by electron transfer between donor and acceptor molecules in contact. The coupling for that transfer depends on the overlap of the two π systems in the same way as the coupling between like molecules, so the relative orientation of donor and acceptor sets the balance between separation and recombination. At the heterointerface itself, face-on orientation of the donor relative to the acceptor raises the yield of separated charge, and the performance of polymer solar cells follows the orientation of the molecules at that interface more closely than the orientation in either bulk layer.\autocite{tumbleston2014influence}

In a light-emitting diode the orientation of the emitting dipoles sets the fraction of the generated light that leaves the device (\cref{fig:2.1}c). Light generated in the organic layers, whose refractive index exceeds those of the glass substrate and of air, escapes only within a narrow cone about the film normal, and light outside that cone is totally reflected at the substrate and air interfaces. A dipole radiates most strongly perpendicular to its own axis, so an emitter whose dipole lies in the plane of the film sends its light toward the viewer, while an emitter whose dipole stands along the film normal sends its light along the film, where it remains trapped in waveguided and surface plasmon modes. For an emitter whose transition dipole lies in the molecular plane, face-on molecules contribute only horizontal dipoles and edge-on molecules contribute vertical components as well, so the fraction of light that escapes follows the distribution of molecular orientation in the emitting layer. Such a preferred orientation can form even in the amorphous films used as emitting layers, which have no crystalline order. Horizontal alignment of the emitters raises the fraction of light that escapes by roughly half relative to an isotropic film.\autocite{yokoyama2011molecular,brutting2013device}

In an organic spin valve the relevant property is the electronic state of the molecules in contact with a ferromagnetic electrode (\cref{fig:2.1}d). The device passes current from one ferromagnetic electrode through the molecular layer to a second, and its resistance changes when the relative magnetization of the two electrodes is reversed, provided the carriers retain a net spin polarization across the layer. Molecules in that first layer hybridize with the metal, which shifts and broadens their frontier orbitals and can polarize them in spin, so the interface acts as a spin filter with properties of its own. The majority and minority spin states of the ferromagnet couple to the molecular orbitals with different strengths, so the interfacial layer can enhance the spin polarization of the injected current or reverse its sign relative to that of the bare electrode.\autocite{sanvito2010rise,barraud2010unravelling} The strength of this coupling depends on the overlap between the molecular orbitals and the metal states, which is largest when the π system lies flat against the electrode and decreases with the height at which the molecule sits, so the spin filtering of the interface follows the orientation and adsorption height of its first molecular layer. The same interaction sets the alignment of energy levels at every metal-organic contact. Charge transfer between molecule and metal and the redistribution of metal electrons beneath the adsorbed layer form an interface dipole that shifts the molecular levels away from the positions predicted by aligning the vacuum levels of the two materials, and with them the barrier to charge injection in each of the devices above.\autocite{ishii1999energy}

Orientation specifies how a rigid molecule sits, yet a molecule at an interface also has internal degrees of freedom, and the conformation it adopts there need not match the one it adopts in the bulk. Torsions between linked rings, the arrangement of ligands about a metal center, and the planarity of a macrocycle all respond to the forces at an interface. A molecule bound to a reactive surface can bend or twist to bring more of its π system into contact, a molecule confined between a substrate and the layer above can be flattened, and a metal complex whose central atom meets the substrate on one face can shift its ligands to accommodate the contact. Even a rigid macrocycle such as ZnPc can in principle bend out of plane when its central ion or its outer rings bind to a substrate, and the quinolinolate ligands of Alq\textsubscript{3}, described in the next section, have more freedom still. Each change alters the overlap of the π system with itself and with the material beneath it, and with it the energies of the frontier orbitals, the strength of hybridization with the substrate, and the directions of the transition dipoles that govern absorption and emission.

Changes of this kind could give an interface properties found in neither the substrate nor the bulk film. A torsion that interrupts the conjugation of a molecule would localize its frontier orbitals and weaken its coupling to neighbors, a distortion that lowers the symmetry of a molecule would split degenerate orbitals and allow transitions that are forbidden in the free molecule, and a change in the bonding at a metal center would alter both the injection barrier and the spin-dependent hybridization described above. Whether such states form at the buried interfaces of working films, and how strongly they matter, has not been established, because no available measurement resolves the conformation of a single molecular layer beneath a film. Diffraction records the periodic packing of crystalline grains averaged through the film, photoemission and scanning probes reach exposed surfaces and submonolayer coverages, which continued deposition buries and can alter, and spectroscopies of the whole film average the interfacial layer with the many layers above it. Interfacial conformation therefore remains an open variable, likely minor in many systems but ruled out in none by present measurements. Its signature would appear in the resonances of the interfacial layer, since a change of conformation shifts the energies, strengths, and dipole directions of the transitions it involves, and the depth-resolved refinement of resonance amplitudes developed in \cref{chap:resonances} is sensitive to such changes even where it cannot identify their cause.

In all four devices the property of interest is set where molecules meet another material, a dielectric, an acceptor, an electrode, or a ferromagnet. At each of these interfaces the orientation, packing, and possibly the conformation of the first molecular layer set the coupling that matters, between molecules in a transistor channel, between donor and acceptor in a solar cell, between emitting dipoles and the optical modes of a diode, and between molecular orbitals and metal states at an electrode. The arrangement and electronic state of thosemolecules are also the quantities that conventional measurements reach least well, since they lie beneath the film and occupy one or two molecular layers. Resolving them requires a probe that reaches a layer of that thickness beneath a finished film and responds to both the orientation and the electronic structure of its molecules. The resonant optical response developed in the rest of this chapter is built for that purpose.

\begin{figure}
  \centering
  \includegraphics[width=\linewidth,height=0.75\textheight,keepaspectratio]{molecular-chapter/molch_growth}
  \caption[Growth modes of a vapor-deposited molecular film.]{Growth modes of a vapor-deposited molecular film. Standing molecules are drawn as slabs on a substrate at nominal coverages of 0.5, 1, and 3 monolayers (ML), with the molecules in contact with the substrate in red. (a) Layer-by-layer (Frank-van der Merwe) growth, in which each layer completes before the next nucleates. (b) Layer-plus-island (Stranski-Krastanov) growth, in which a complete wetting layer forms before three-dimensional islands nucleate on it. (c) Island (Volmer-Weber) growth, in which islands nucleate directly on the substrate and leave it partly exposed until they coalesce.}
  \label{fig:2.2}
\end{figure}

\section{Growth of Vapor-Deposited Molecular Films}\label{sec:growth-of-vapor-deposited-molecular-films}

The films studied here are grown by physical vapor deposition, in which molecules sublimed from a heated source condense on a substrate held in vacuum. The arrangement they adopt is set during growth by a competition between the attraction of each arriving molecule to the substrate and its attraction to the molecules already deposited.\autocite{venables1984nucleation,bauer1958phanomenologische} Three growth modes bound the possibilities (\cref{fig:2.2}). When adhesion to the substrate dominates, molecules spread into a complete layer before the next layer nucleates, the layer-by-layer or Frank-van der Merwe mode. When cohesion dominates, molecules gather into three-dimensional islands that leave the substrate exposed until the islands coalesce, the island or Volmer-Weber mode. Layer-plus-island or Stranski-Krastanov growth lies between the two, with a wetting layer that forms first and islands that grow on it.

Molecular films depart from this atomic picture in two ways that matter for the interface. A planar molecule binds to a substrate and to its neighbors with energies that depend on its orientation, so the competition between adhesion and cohesion selects an orientation as well as a growth mode.\autocite{forrest1997ultrathin,witte2004growth} Molecules lying flat maximize their contact with a strongly interacting substrate, molecules standing on edge maximize their π contact with one another, and the orientation adopted in the first layer need not persist once that layer is covered. Growth is also kinetic. The deposition rate and the substrate temperature set how far each molecule diffuses before it is buried, so the structure of a film can change with thickness, and its first layers can be trapped in states that would not form at equilibrium.\autocite{kowarik2008organic}

Zinc phthalocyanine (ZnPc) is the template for this discussion and for the measurements of \cref{chap:znpc,chap:resonances}. It is a planar, rigid macrocycle of \(D_{4h}\) symmetry in which four isoindole units surround a central zinc ion, and its carbon atoms fall into four inequivalent sites, each with its own core level (\cref{fig:4.1}). On native-oxide silicon ZnPc forms polycrystalline films in a herringbone packing with the molecules standing edge-on, while on templating layers such as CuI and graphene it adopts a brickstone packing with the molecules lying face-on.\autocite{kato2020morphological,ngongangndjawa2015impact,schunemann2011zinc} The templating picture assigns the orientation of the whole film to the first layer, much as a seed layer sets the orientation of an epitaxial inorganic film. Coverage-checkpoint studies, which interrupt deposition to measure a partial film, test this picture and disagree with one another. In situ measurements resolve a density stratification inconsistent with simple templating\autocite{derrien2022interfacial}, while other work argues that an exposed growth front records a transient state that changes once it is buried.\autocite{yamamoto2022variation,deoteyza2006structural}

Tris(8-hydroxyquinolinato)aluminum (Alq\textsubscript{3}), the second molecule studied here, shows why the electronic state of a molecule in a film cannot be assumed from its formula. Three quinolinolate ligands coordinate a central aluminum ion, each through one oxygen and one nitrogen, and the six donor atoms can take two arrangements. The meridional (mer) isomer places the three oxygens in one plane through the metal, and the facial (fac) isomer places them on one face of the coordination octahedron.\autocite{brinkmann2000correlation,colle2003preparation} Vapor-deposited films consist mainly of the mer isomer, are amorphous with weak orientational order, and leave the ligands some conformational freedom about the metal. Alq\textsubscript{3} was the emitter and electron transporter of the first efficient organic light-emitting diode\autocite{tang1987organic} and remains a standard molecule in studies of organic spin valves.\autocite{xiong2004giant,barraud2010unravelling} Its role in this work complements that of ZnPc. ZnPc tests whether orientation can be resolved at a buried interface, and Alq\textsubscript{3} tests whether the electronic state of a molecule, as recorded in its NEXAFS, is the same in every film. \Cref{chap:resonances} takes up that question without assigning the differences it finds to a particular cause.

Films of both molecules leave open where the first molecular layer sits and why, whether its orientation and density follow the substrate, follow the bulk, or take a form of their own. They also leave open whether molecules within a monolayer of an interface occupy electronic states that differ from those of the bulk, through hybridization with the substrate, a change of conformation, or a change of packing. Measurements of an exposed growth front cannot settle either problem, since continued deposition buries the interface and can alter it. Both require a measurement of orientation and electronic state at the buried interface of a finished film, which the rest of this chapter prepares optically and \cref{chap:znpc,chap:resonances} carry out.

\section{Core Excitations and Transition Dipoles}\label{sec:core-excitations-and-transition-dipoles}

\begin{figure}[htp]
  \centering
  \includegraphics[width=\linewidth,height=0.75\textheight,keepaspectratio]{molecular-chapter/molch_orbitals}
  \caption[Orbitals and core transitions of a conjugated carbon bond.]{Orbitals and core transitions of a conjugated carbon bond. Energies are not to scale. (a) Mixing of the carbon 2s orbital with two 2p orbitals gives three sp2 hybrids, while the remaining pz orbital is unchanged. (b) Hybrids on two neighboring carbons form the σ and σ* orbitals of the C=C bond, and their pz orbitals form π and π*. Absorption of an X-ray photon promotes a C 1s electron into the empty π* (red) or σ* (gray) level. (c) Shapes of the four orbitals along the bond, with filled and open lobes of opposite phase, bonding orbitals below, antibonding orbitals above, and dashed lines at nodal planes. The transition dipole μ of C 1s → π* lies along the p orbitals, normal to the molecular plane, while the dipole of C 1s → σ* lies along the bond.}
  \label{fig:2.3}
\end{figure}

The resonances that give soft X-ray measurements their chemical and orientational sensitivity are excitations of a core electron into unoccupied molecular orbitals, and in a conjugated molecule those orbitals follow from the hybridization of carbon (\cref{fig:2.3}). The 2s orbital and two of the 2p orbitals of each carbon mix into three sp\(^{2}\) hybrids that lie in the molecular plane at \ang{120} to one another, and the remaining 2p orbital, \(p_{z}\), stands normal to that plane (\cref{fig:2.3}a). Hybrids on neighboring atoms overlap head-on to form the σ bonds of the molecular framework, and the \(p_{z}\) orbitals overlap side-on to form the delocalized π system (\cref{fig:2.3}b). Each overlap yields a bonding orbital, occupied in the ground state, and an antibonding orbital, empty in the ground state and available as the final state of an X-ray excitation.

Side-on overlap is weaker than head-on overlap, so the π pair is split less than the σ pair and the π* orbitals lie lowest among the empty states. The carbon K edge shows this ordering directly. Excitations from C 1s into π* orbitals give sharp resonances between about 284 and \qty{290}{eV}, and excitations into σ* orbitals, most of which lie above the ionization threshold and couple to the continuum, give broader resonances from about \qty{290}{eV} upward.\autocite{stohr2010nexafs} The C 1s orbital is confined to one atom and barely perturbed by bonding, so each excitation starts from a known site, and carbons in different chemical environments contribute resonances at different energies.

The absorption spectra are set by the coherent sum of electronic transitions excited by incident photons. For X-ray photons incident on molecular materials, the initial state is commonly an atomic orbital deep within the atomic potential well and minimally perturbed by the molecular bonds (\cref{fig:2.3}b). The final states, however, are delocalized unoccupied molecular orbitals that commonly stretch over a large portion of the molecule. Fermi's golden rule relates the transition dipole moment (TDM) between the initial and final states to the transition probability, i.e.,

\begin{equation}
  \Gamma_{ij} \propto \bigl|\langle f \,|\, \mathbf{\mu} \cdot \mathbf{E} \,|\, i \rangle\bigr|^{2},
  \label{eq:golden-rule}
\end{equation}

where \(\mathbf{\mu}\) is the dipole operator and \(\mathbf{E}\) the electric field of the incident photon. The oscillator strength of the transition, \(f_{if} \propto \omega_{fi}\,|\langle f \,|\, \widehat{\mathbf{e}} \cdot \mathbf{\mu} \,|\, i \rangle|^{2}\) with \(\widehat{\mathbf{e}}\) the unit polarization vector, sets the integrated absorption cross section of the corresponding resonance and therefore its area in the absorption spectrum.\autocite{stohr2010nexafs} \Cref{chap:resonances} tracks this quantity for individual resonances. Because the transition probability depends on the alignment between the polarization and the transition dipole, the absorption is a tensor quantity. In scattering this appears as a polarization-dependent contrast tied to molecular orientation and to the bonding conformations that set the relative orientation of the transition dipoles within the film.

The direction of each transition dipole follows from the shape of the final orbital (\cref{fig:2.3}c). The C 1s orbital is spherical, so the dipole operator connects it only to the p character of the final state on the excited atom, and the transition dipole points along that p component. For a π* final state, built from \(p_{z}\) orbitals, the dipole lies along the normal to the molecular plane. For a σ* final state, built from hybrids directed along a bond, the dipole lies along that bond and therefore in the molecular plane. Every resonance of a planar conjugated molecule thus carries a direction fixed in the frame of the molecule, and its absorption scales as the squared projection of the electric field on that direction.

\section{Chemical and Orientational Contrast}\label{sec:chemical-and-orientational-contrast}

\begin{figure}[htp]
  \centering
  \includegraphics[width=0.6\linewidth,height=0.75\textheight,keepaspectratio]{fig2-04_polymer-nexafs}
  \caption[NEXAFS spectroscopy measurements comparing polymers.]{NEXAFS spectroscopy measurements comparing polymers. Polystyrene (red), poly(methyl methacrylate) (blue), and poly(vinyl pyridine) (green), with the repeat unit of each drawn in the box of matching color.}
  \label{fig:2.4}
\end{figure}

The same transitions give soft X-ray measurements their chemical contrast (\cref{fig:2.4}). Polystyrene and poly(vinyl pyridine) carry aromatic rings and show a strong C 1s → π* resonance of the ring near \qty{285}{eV}, split in poly(vinyl pyridine) by the nitrogen in its ring. Poly(methyl methacrylate) has no ring, and its strongest π* resonance, near \qty{288.5}{eV}, comes from the carbonyl of its ester side group. Three polymers built from the same few elements are therefore distinguished by the energies of their resonances alone, and resonant scattering converts this distinction into contrast between materials.\autocite{ade2008nexafs}

Orientation adds a second source of contrast. Each resonance absorbs according to the projection of the field on its transition dipole, so two layers of the same molecule have different optical constants when their molecules are oriented differently, and the difference is largest at the energies of the most strongly aligned resonances. Polarized resonant scattering uses this to distinguish regions that are chemically identical and differ only in molecular orientation.\autocite{collins2012polarized} The same sensitivity means that a molecular film cannot be described by a single complex index. Its optical response is a tensor whose components follow from the dipoles of its transitions and the orientation of its molecules, and the remainder of this chapter constructs that tensor, beginning with the oscillator of a single transition.

\section{The Resonant Oscillator and Causality}\label{sec:the-resonant-oscillator-and-causality}

A transition dipole fixes how strongly and along which direction a resonance absorbs. The optical response of a material also contains the dispersion that accompanies each absorption, and both follow from treating each transition as a driven, damped oscillator. The treatment is semiclassical. Quantum mechanics supplies the energy, oscillator strength, and dipole of each transition through Fermi's golden rule, and the classical response of a bound charge supplies the line shape and its causal structure.\autocite{alsnielsen2011elements,jackson2009classical,kuzyk2023nonlinear} An electron bound with resonant frequency \(\omega_{j}\) and damping \(\Gamma_{j}\), driven by a field \(\mathbf{E}_{0}e^{- i\omega t}\), obeys

\begin{equation}
  m\left( \ddot{x} + \Gamma_{j}\dot{x} + \omega_{j}^{2}x \right) = - e\,{\widehat{\mathbf{\mu}}}_{j} \cdot \mathbf{E}_{0}\, e^{- i\omega t},
  \label{eq:oscillator}
\end{equation}

where \(x\) is the displacement along the unit transition dipole \({\widehat{\mathbf{\mu}}}_{j}\), so only the projection of the field on the dipole drives the oscillator. The steady-state displacement induces a dipole along \({\widehat{\mathbf{\mu}}}_{j}\), and summing over transitions of oscillator strength \(f_{j}\) in a material of \(N\) molecules per unit volume gives the susceptibility tensor

\begin{equation}
  \chi_{ab}(\omega) = \frac{Ne^{2}}{\varepsilon_{0}m}\sum_{j}\frac{f_{j}\,{\widehat{\mu}}_{j,a}\,{\widehat{\mu}}_{j,b}}{\omega_{j}^{2} - \omega^{2} - i\Gamma_{j}\omega}.
  \label{eq:susceptibility}
\end{equation}

Each transition contributes a tensor of rank one along its own dipole. For light polarized along \(\widehat{\mathbf{e}}\) the absorption is proportional to the imaginary part of \(\widehat{\mathbf{e}} \cdot \mathbf{\chi} \cdot \widehat{\mathbf{e}}\), a sum over transitions weighted by \(\left( \widehat{\mathbf{e}} \cdot {\widehat{\mathbf{\mu}}}_{j} \right)^{2}\), which is the polarization dependence given by the golden rule. Near a resonance the imaginary part is a Lorentzian of width \(\Gamma_{j}\), and the real part changes sign across it, positive below the resonance and negative above.

The pairing of absorption and dispersion holds beyond the oscillator model. A physical response cannot precede its cause, so the response function \(\chi(t)\) vanishes for \(t < 0\) and its Fourier transform \(\chi(\omega)\) is analytic in the upper half of the complex frequency plane.\autocite{toll1956causality} Cauchy's theorem on a contour closed in that half plane ties the real and imaginary parts of \(\chi\) together, and with the symmetry \(\chi( - \omega) = \chi^{*}(\omega)\) of a real response the relation becomes

\begin{equation}
  \operatorname{Re}\chi(\omega) = \frac{2}{\pi}\,\mathcal{P}\int_{0}^{\infty}\frac{\omega'\, \operatorname{Im}\chi(\omega')}{\omega'^{2} - \omega^{2}}\, d\omega'.
  \label{eq:kk-chi}
\end{equation}

These are the Kramers-Kronig relations.\autocite{kramers1927diffusion,delkronig1926theory} They require no model of the individual transitions and hold for each component of the tensor separately. The index of refraction tensor is the standard descriptor of the optical response within the soft matter X-ray community, written here as \(n = 1 - \delta + i\beta\) for convenience with \(\beta\) directly corresponding with the absorption, and \(\delta\) describing the phase. Written for the refractive index and the photon energy, the same relation becomes the transform of \(\beta\) into \(\delta\) given below. The atomic scattering factors \(f_{1}\ \&\ f_{2}\), the permittivity \(\epsilon\), and the susceptibility \(\chi\) follow from the same response as

\begin{align}
  \delta &= \frac{r_{e}\lambda^{2}}{2\pi}n_{a}f_{1}, & \beta &= \frac{r_{e}\lambda^{2}}{2\pi}n_{a}f_{2}, \label{eq:delta-beta}\\
  \epsilon &= n^{2} \simeq 1 - 2\delta + 2i\beta, & \chi &= \epsilon - 1 \simeq - 2\delta + 2i\beta, \label{eq:epsilon-chi}
\end{align}

where \(r_{e}\) is the classical electron radius and \(n_{a}\) the number density of atoms.\autocite{attwood1999soft} The optical response manifests from the ensemble of electronic transitions excited by incident photons, resulting in an energy dependent, dispersive optical response. The energy levels and transition intensities are set by the local bonding configurations, heteroatoms, and molecular conformations, so the resonant structure contains the chemical state. Through the index contrast, this chemical state appears in the scattering profile, presenting the pathway to resolve chemical compositions within these scattering structures.

The resonant optical response is most directly measured through the mass absorption \(\mu_{m}(E)\), which the Beer-Lambert law relates to the absorptive part of the index,

\begin{equation}
  \beta(E) = \frac{\lambda\rho\,\mu_{m}(E)}{4\pi},
  \label{eq:beta-mu}
\end{equation}

where \(\rho\) is the mass density and \(\lambda = hc/E\) the wavelength. The real component is then computed using the Kramers-Kronig relationship that causally constrains the optical response. The relationship is equivalent to a scaled Hilbert transform, and requires the Cauchy principal value integral over the total energy domain

\begin{equation}
  \delta(E) = \frac{2}{\pi}\,\mathcal{P}\int_{0}^{\infty}\frac{E'\beta(E')}{E^{2} - E'^{2}}\, dE'.
  \label{eq:kk-delta}
\end{equation}

At every energy, \(\delta\) depends on the absorption of every other energy, so the transformation requires extrapolation to 0 and \(\infty\) respectively. This is most often achieved with tabulated bare-atom spectra, which are accurate far from chemical resonance.\autocite{henke1993x} This need for extrapolation presents challenges to the method that are addressed in the following chapters.

The line shape of the oscillator is Lorentzian, set by the lifetime of the core hole, which broadens each transition at the carbon K edge by about \qty{0.1}{eV}. Resonances measured in molecular films are broader and closer to Gaussian, because vibrational structure, inequivalent environments, and the resolution of the monochromator each spread a transition over a distribution of energies. A Gaussian absorption describes such an ensemble, and its causal partner follows from the same Kramers-Kronig relation. For a transition of amplitude \(A_{j}\) centered at \(E_{j}\) with width \(\sigma_{j}\),

\begin{equation}
  \begin{gathered}
  \beta_{j}(E) = A_{j}\exp\left[ - \frac{\left( E - E_{j} \right)^{2}}{2\sigma_{j}^{2}} \right], \\
  \delta_{j}(E) \simeq \frac{2A_{j}}{\sqrt{\pi}}\, D\left( \frac{E - E_{j}}{\sqrt{2}\,\sigma_{j}} \right),
  \end{gathered}
  \label{eq:gauss-dawson}
\end{equation}

where \(D(x) = e^{- x^{2}}\int_{0}^{x}e^{t^{2}}\, dt\) is the Dawson function. The approximation keeps the resonance at \(+ E_{j}\) and drops its mirror at \(- E_{j}\), which the X-ray form of the relation retains through the \(E^{2} - E'^{2}\) in its denominator. Near \qty{285}{eV} the mirror lies about \qty{570}{eV} away and adds only a small, slowly varying offset across the edge, so each Gaussian and Dawson pair is causal to well within the precision of the measurements made here. \Cref{chap:resonances} writes the full optical response of a film as a sum of such pairs.

\begin{figure}
  \centering
  \includegraphics[width=\linewidth,height=0.75\textheight,keepaspectratio]{fig2-03_optical-constants}
  \caption[Optical constants of the \cref{chap:znpc} ZnPc film.]{Optical constants of the \cref{chap:znpc} ZnPc film. (a) In-plane \(\beta_{xx}\) (red) and out-of-plane \(\beta_{zz}\) (black) from the refined building-block model, with the bare-atom absorption (gray, dashed) and the pre-edge (270 to \qty{282}{eV}) and post-edge (320 to \qty{330}{eV}) normalization windows shaded. (b) \(\delta_{xx}\) and \(\delta_{zz}\) from the Kramers-Kronig transform of β. The inset shows the bare-atom β from \qty{10}{eV} to \qty{30}{keV} that closes the transform, with the measured window in red. (c) Relations between the refractive index n, the atomic scattering factors f1 and f2, the permittivity ε, and the susceptibility χ, with β = λρμm/(4π) linking β to the mass absorption μm.}
  \label{fig:2.5}
\end{figure}

\Cref{fig:2.5} carries these relations through for the ZnPc film of \cref{chap:znpc}. The absorptive part follows from the mass absorption through \(\beta = \lambda\rho\mu_{m}/(4\pi)\),\autocite{attwood1999soft} with the measured spectrum scaled to the bare-atom absorption of ZnPc\autocite{henke1993x} in a pre-edge window (270 to \qty{282}{eV}) and a post-edge window (320 to \qty{330}{eV}), where molecular bonding no longer alters the cross section (\cref{fig:2.5}a). Between the windows the edge-on texture splits \(\beta\) into two tensor components. The \(\pi^{*}\) resonances near \qty{285}{eV} appear almost entirely in \(\beta_{xx}\), since their transition dipoles lie along the ring normal and therefore close to the film plane, while the \(\sigma^{*}\) band above \qty{290}{eV} is stronger in \(\beta_{zz}\). The dispersive part follows from the Kramers-Kronig relation above (\cref{fig:2.5}b), evaluated by exact integration of a piecewise-linear \(\beta\), a special case of the piecewise-polynomial method of Watts.\autocite{watts2014calculation} Because \(\delta\) at each energy depends on \(\beta\) at every energy, the measured window is embedded in the bare-atom \(\beta\) from \qty{10}{eV} to \qty{30}{keV} before the transform is evaluated (\cref{fig:2.5}b, inset). The transform turns each absorption peak into a dispersive feature, so \(\delta_{xx}\) reaches about \(- 3 \times 10^{- 3}\) just below the first \(\pi^{*}\) resonance and changes sign across it, and both components approach a positive, slowly varying value well above the edge. Contrast in a resonant measurement therefore changes in sign as well as magnitude with energy. The energy of greatest contrast between two components therefore need not coincide with an absorption peak. \Cref{fig:2.5}c collects the equivalent descriptions of the same response, the atomic scattering factors \(f_{1}\) and \(f_{2}\), the permittivity \(\epsilon = n^{2}\), and the susceptibility \(\chi = \epsilon - 1\), so that results stated in any of these conventions can be compared directly.\autocite{attwood1999soft,yan2013accurate}

A fit that treats δ and β as independent at every energy can return optical constants that no causal material could have, with the dispersion at one energy unrelated to the absorption across the edge. Tying every resonance to a Gaussian and Dawson pair removes that freedom, since the dispersion at any energythen follows from the same few parameters that describe the absorption. Those parameters are the energy, width, strength, and dipole of each transition, and the next section takes them from DFT.

\section{Transitions from ML-Refined DFT}\label{sec:transitions-from-ml-refined-dft}

\begin{figure}
  \centering
  \includegraphics[width=\linewidth,height=0.75\textheight,keepaspectratio]{molecular-chapter/molch_dft}
  \caption[From DFT transitions to a refined cluster model of ZnPc.]{From DFT transitions to a refined cluster model of ZnPc. (a) Spectrum broadened from the 104 DFT transitions (black) against the sum of the 19 cluster Gaussians (red, dashed), on the DFT energy scale and normalized to its maximum. Below, each transition as a stick of height proportional to its oscillator strength, alternating in color between adjacent clusters. (b) Angle-resolved NEXAFS at θ = \ang{20} to \ang{90} (points) with the refined cluster model (lines), and the residuals as a percentage of the largest measured intensity. (c) Change of each cluster parameter from the DFT seed to the refined model, by cluster index in order of energy. From top, the change in full width at half maximum, the transition dipole angle ϑ from the ring normal for the seed (gray) and the refined model (red), its change Δϑ, the energy shift normalized by the seed width, and the logarithm of the ratio of refined to seed amplitude, each set of amplitudes normalized to its total. Seed energies include the rigid shift of \qty{0.45}{eV}.}
  \label{fig:2.6}
\end{figure}

The oscillators of a molecule can be taken from a measured spectrum or computed from its electronic structure. A measured spectrum fixes the total absorption at each energy but cannot say which transitions produce it, along which directions their dipoles point, or how much of the measured dichroism belongs to the molecule rather than to its orientation in the film. A calculation supplies all three, but its absolute energies carry an offset, its line widths are imposed by a chosen broadening rather than computed, and it treats an isolated molecule without the intermolecular interactions or vibronic structure of a film, so its spectrum departs from a measured one in energy, width, and relative strength. The approach used here, developed by Murcia,\autocite{murcia2026quantitative} combines the two (\cref{fig:2.6}). DFT supplies the transitions, clustering reduces them to the few that a spectrum can resolve, and refinement against angle-resolved NEXAFS adjusts the clusters while keeping the information the calculation carried.

The transitions of ZnPc were computed for a single molecule with the StoBe code by the transition-potential method with a half core hole, which yields the energy, oscillator strength, and transition dipole of every excitation, while \(\Delta KS\) calculations of the ground and core-excited states fix the absolute energy of each excitation center.\autocite{murcia2026quantitative} One calculation is made for each of the four inequivalent carbon sites, with the molecular geometry taken from the crystal structure, which agreed with experiment better than geometries relaxed thermodynamically or with force fields (Supplemental Material). Each transition carries the site from which it starts and the molecular orbital into which it excites, so every feature of the computed spectrum can be traced to a bond. The calculation treats an isolated molecule and so contains no states shared between neighbors, and its absolute energy scale carries an error of a few tenths of an electron volt, corrected by a rigid shift of \qty{0.45}{eV} before comparison with experiment.

A calculation of this kind returns more transitions than any spectrum resolves. The filtered calculation for ZnPc contains 104 transitions between 283.8 and \qty{318.4}{eV}, many of them separated by less than their own width and sharing a dipole direction (\cref{fig:2.6}a). Hierarchical agglomerative clustering groups such transitions into clusters, each standing for a family that no measurement could separate and carrying its own energy, width, strength, and dipole. Two hyperparameters control the grouping, a threshold on oscillator strength (OS) and a threshold on the overlap between neighboring line shapes (OVP). Both were scanned and each clustering compared with measured NEXAFS and with the original calculation, which selected OS = 0.25 and OVP = 52.5, with similar results for any OVP between 50 and 55 (Supplemental Material). The resulting 19 clusters reproduce the broadened calculation across the edge, with small departures near \qty{290}{eV} and at the highest energies.

Refinement then fits the clusters to angle-resolved NEXAFS of an edge-on film measured at five angles (\cref{fig:2.6}b). The energy, width, strength, and dipole direction of each cluster are adjusted together with the tilt that projects the molecular tensor into the frame of the film. The refined model reproduces the spectra at every angle with an average deviation of 2.7 percent (\cref{chap:znpc}), and its largest residuals fall on the narrow π* features. The changes the refinement makes are recorded cluster by cluster in \cref{fig:2.6}c. Widths broaden by up to about a quarter of an electron volt except in the two highest clusters, which narrow. Energies move by less than a width. The dipoles of the lowest π* clusters stay along the ring normal, the σ* clusters rotate further into the ring plane, and strength moves from the broad clusters above the edge toward the resolved resonances below it. Because each change is modest, every refined cluster remains identifiable with its seed and keeps the site and orbital assignment inherited from DFT.

The construction has two limits, both anticipated in the first question of \cref{chap:introduction}. The refinement is made against a reference film whose own orientation enters its dichroism, so the dipole of each cluster and the tilt of the film are fixed by the same data and are correlated. Clustering also keeps only the transitions selected as principal, so a weak transition, or one that appears only in a different molecular environment, has no cluster to represent it. \Cref{chap:znpc} tests the first limit by comparing the tilt recovered from reflectivity with the tilt of the reference, and \cref{chap:resonances} tests the second by refining the strength of individual clusters with depth.

\section{Molecular and Film Reference Frames}\label{sec:molecular-and-film-reference-frames}

The refined clusters describe one molecule. Each cluster is a Gaussian and Dawson pair attached to a unit dipole, and together with the bare-atom absorption that carries the edge step they fix the response of the molecule at every energy and polarization. A film contains many molecules in many orientations, and its response follows from how the frame of the molecule is distributed relative to the substrate.

The anisotropy of a molecular film originates in the orientation of its transition dipole moments. Each resonance couples most strongly to light polarized along its transition dipole and not at all to light polarized perpendicular to it, so the molecular optical tensor is a sum over transitions of the outer product \(\widehat{\mu}{\widehat{\mu}}^{T}\) of each unit transition dipole, weighted by the oscillator strength and line shape of that transition. For a planar conjugated molecule the C 1s → π* dipoles lie along the ring normal while the σ* dipoles lie largely in the ring plane, so the π* resonances enter the out-of-plane component of the molecular tensor and the σ* resonances its in-plane components (\cref{fig:2.3}c). The film tensor then follows from how the molecular frame is distributed relative to the substrate, and takes one of the three forms below.

\begin{equation}
  n_{\mathrm{iso}} = \begin{pmatrix} xx & 0 & 0 \\ 0 & xx & 0 \\ 0 & 0 & xx \end{pmatrix},\quad
  n_{\mathrm{uni}} = \begin{pmatrix} xx & 0 & 0 \\ 0 & xx & 0 \\ 0 & 0 & zz \end{pmatrix},\quad
  n_{\mathrm{bi}} = \begin{pmatrix} xx & 0 & 0 \\ 0 & yy & 0 \\ 0 & 0 & zz \end{pmatrix}.
  \label{eq:tensor-symmetry}
\end{equation}

The most general film response is a fully biaxial tensor, but most molecular films instead adopt one of the two simpler limits shown above. An isotropic film, with all three axes equal, requires the molecular ensemble to be fully randomized in three dimensions. More commonly due to processing methods, physical vapor deposition and solvent process techniques like spin coating order a single molecular axis, such as the macrocycle normal, relative to the substrate normal while leaving the in-plane azimuthal distribution unconstrained, producing the uniaxial case considered throughout this work.

\begin{figure}[htp]
  \centering
  \includegraphics[width=0.75\linewidth,height=0.75\textheight,keepaspectratio]{fig2-07_tensor-construction}
  \caption[Schematic construction of a uniaxial molecular film.]{Schematic construction of a uniaxial molecular film. Molecular crystallites are drawn with their transition dipole moments μ. (a) In-plane distribution of crystallites, evenly distributed in azimuth across the film. (b) Film cross section with distinct in-plane and out-of-plane averaged dipoles. (c) Ensemble average over the tilt γ between the molecular plane and the substrate, equal to the tilt of the ring normal from the surface normal, and over in-plane rotations ϑ.}
  \label{fig:2.7}
\end{figure}

Explicitly, a uniaxial molecular tensor \(n^{m} = \operatorname{diag}\left( n_{xx}^{m},n_{xx}^{m},n_{zz}^{m} \right)\) tilted by \(\gamma\) about the in-plane \(x\) axis and averaged over the azimuthal rotation \(\vartheta\) sketched in \cref{fig:2.7} gives the film tensor

\begin{equation}
  n^{f}(E;\gamma,\rho) = \frac{\rho}{2\pi}\int_{0}^{2\pi} R_{z}(\vartheta)\, R_{x}(\gamma)\, n^{m}(E)\, R_{x}(\gamma)^{\mathsf{T}}\, R_{z}(\vartheta)^{\mathsf{T}}\, d\vartheta,
  \label{eq:film-tensor}
\end{equation}

whose diagonal components are

\begin{align}
  n_{xx}^{f}(E;\gamma,\rho) &= \frac{\rho}{2}\left[ n_{xx}^{m}\left( 1 + \cos^{2}\gamma \right) + n_{zz}^{m}\sin^{2}\gamma \right], \label{eq:film-xx}\\
  n_{zz}^{f}(E;\gamma,\rho) &= \rho\left( n_{xx}^{m}\sin^{2}\gamma + n_{zz}^{m}\cos^{2}\gamma \right), \label{eq:film-zz}
\end{align}

recovering the molecular values at \(\gamma\  = \ 0\) as expected for a rotation of a uniaxial tensor. \(n_{xx}\) carries an additional factor of \((1 + cos²\gamma)/2\) rather than simply \(cos²\gamma\) because it is further averaged over the in-plane azimuthal angle \(\vartheta\), whereas \(n_{zz}\), fixed along the rotation axis, requires no such averaging. The trace is conserved, \(2n_{xx}^{f} + n_{zz}^{f} = \rho\left( 2n_{xx}^{m} + n_{zz}^{m} \right)\) for every \(\gamma\), which provides a direct check of any implementation. Supplemental Material extends this relation to a partially ordered ensemble of finite density \emph{ρ}, the form used throughout the structural fits in \cref{chap:znpc,chap:resonances}.

\section{Angle-Resolved NEXAFS of a Uniform Film}\label{sec:angle-resolved-nexafs-of-a-uniform-film}

\begin{figure}[htp]
  \centering
  \includegraphics[width=\linewidth,height=0.75\textheight,keepaspectratio]{fig1-02_ar-nexafs}
  \caption[Angle-resolved NEXAFS of an edge-on ZnPc film, resolved by transition cluster.]{Angle-resolved NEXAFS of an edge-on ZnPc film, resolved by transition cluster. (a) Energy diagram for a C 1s → π* excitation, and a conjugated carbon ring above the substrate with the p orbitals that form its π* states normal to the ring plane. The transition dipole μ of a C 1s → π* excitation points along the ring normal, at angle α from the surface normal n. (b) Mass absorption \(\mu_{m}\) from 283.5 to \qty{300}{eV} with the electric field at θ = \ang{20} to \ang{90} from the surface normal (light to dark gray points), with fits built from the cDFT clusters (lines). Triangles mark the six clusters shown in c. (c) Amplitude against θ for six representative clusters, normalized by the isotropic amplitude \(A_{0}\), with fits to the vector-orbital form \(A_{0}[1 + (1/2)(3\cos^{2}\theta - 1)(3\cos^{2}\alpha - 1)]\). The fitted transition dipole angle α is printed at the right of each curve.}
  \label{fig:2.8}
\end{figure}

Angle-resolved NEXAFS of the edge-on ZnPc film tests the simplest film this construction allows, one in which every molecule shares a single tilt (\cref{fig:2.8})\autocite{stohr2010nexafs,ade2008nexafs}. The \(\pi^{*}\) states are built from p orbitals that stand above and below the conjugated ring, so the transition dipole \(\mathbf{\mu}\) of a C 1s \(\rightarrow \pi^{*}\) excitation points along the ring normal, at an angle \(\alpha\) from the surface normal \(\widehat{\mathbf{n}}\) of a film. By Fermi's golden rule the absorption scales as \(\left| \mathbf{E} \cdot \mathbf{\mu} \right|^{2}\), so rotating \(\mathbf{E}\) from \(\theta\) = \ang{20} to \ang{90} relative to \(\widehat{\mathbf{n}}\) in an edge-on ZnPc film raises the \(\pi^{*}\) resonances near \qty{285}{eV} and lowers the \(\sigma^{*}\) band above \qty{291}{eV} (\cref{fig:2.8}b). Integrating fixed energy windows would mix overlapping transitions under one assumed orbital symmetry. Instead, each spectrum is decomposed into the 16 cDFT transition clusters between 283.5 and \qty{300}{eV}, with centers and widths fixed and one non-negative amplitude per cluster and angle, which reproduces every spectrum with a relative residual below 3 percent. The amplitude of each cluster follows the vector-orbital form \(A_{0}\,[1 + (1/2)(3\cos^{2}\theta - 1)(3\cos^{2}\alpha - 1)]\) (\cref{fig:2.8}c), which gives each cluster its own dipole angle \(\alpha\). The three lowest \(\pi^{*}\) clusters give \(\alpha\) = 74.0 to \ang{74.8}. Clusters between 286 and \qty{291}{eV} fall progressively from \ang{72} to \ang{60}, while the \(\sigma^{*}\) clusters above \qty{293}{eV} give 47 to \ang{49}, below the magic angle of \ang{54.7} at which a uniaxial film shows no dichroism. In a uniaxial film these angles map onto one ring tilt \(\gamma\) through the direction of each dipole in the molecular frame, with \(\alpha = \gamma\) for a dipole along the ring normal and \(\cos^{2}\alpha = (1/2)\sin^{2}\gamma\) for one in the ring plane. A single tilt of \ang{68.2} reproduces the \(\sigma^{*}\) clusters but leaves the \(\pi^{*}\) clusters \ang{6} steeper, a discrepancy that no uniform tilt removes. Angle-resolved NEXAFS therefore measures orientation cluster by cluster, but recovering the molecular state from it requires the tensor treatment developed in this chapter and in \cref{chap:methods,chap:znpc}, which places the bulk tilt at \(\gamma\) = \ang{69.6} within a depth gradient.

Four values of the ZnPc tilt appear in this dissertation, and each answers a different question. The value of \ang{62} is the starting point of the refinement and carries no information about the film. The final step of the refinement fits one tilt to the full set of spectra with the dipoles at their calculated directions and returns \ang{74}, the value reported in \cref{chap:znpc} and Supplemental Material. That fit is dominated by the strong π* resonances, whose dipoles lie along the ring normal and report the tilt directly, and it agrees with the 74.0 to \ang{74.8} found above for the three lowest π* clusters. The single tilt that best reproduces the cluster angles is \ang{68.2}, set mainly by the σ* clusters. Polarized reflectivity resolves the film in depth and places the bulk at \ang{69.6(2)}, with \ang{58.7(6)} at the free surface and \ang{38(1)} at the oxide (\cref{tab:4.1}). A depth gradient of this kind, weighted differently by each measurement, is consistent with the spread between \ang{68} and \ang{74}, although a small error in the calculated σ* dipole directions could contribute as well.

A film built in this way has a fixed number of optical parameters. The refined clusters fix the response of the molecule at every energy and polarization, so each layer of a film is specified by its tilt and density however many energies are measured. This property underlies both questions of the dissertation. It makes the structural fits of \cref{chap:znpc} tractable with an order of magnitude fewer parameters than a free tensor, and it leaves room in \cref{chap:resonances} to refine the strength and orientation of individual resonances with depth without the fit becoming underdetermined.

The failure of a uniform tilt to describe ZnPc motivates resolving the film in depth, which the reflectivity methods of \cref{chap:methods} provide. The construction also rests on a reference spectrum measured once and then reused, which raises the third open problem of this work, whether a spectrum measured on one film can be trusted when it is applied to another. \Cref{chap:xray-atlas} addresses that problem through a shared, versioned record of spectra and of the optical models built from them.
